\documentclass[11pt]{article}
\usepackage{mathblatt}
\makeatletter
\setlength{\mbpfbe}{0mm}% keine Punkte (Bauregel 6.4)
\makeatother
% Rechenraum als Karo (Bauregel 6.9): n Rechenschritte = n mal 10 mm
\newcommand{\karo}[1]{\par\smallskip\noindent\begin{tikzpicture}
  \clip (0,0) rectangle ({\linewidth-1mm},{#1*10mm});
  \draw[black!16,line width=0.3pt,step=5mm] (0,0) grid ({\linewidth},{#1*10mm});
  \end{tikzpicture}\par}
\newcommand{\linie}[1]{\,{\color{black!45}\rule[-1pt]{#1}{0.4pt}}\,}
\newcommand{\pkt}{\textperiodcentered{}}
\begin{document}

\pfheftstil
\pflaufkopf[P10 \pkt{} mehr]{Länge mit Pythagoras \pkt{} Lösungen}{PG50-3}\pfseitevon
\pfheftkopf{Länge mit Pythagoras}{Lösungen \pkt{} mehr}
\begin{pfloesungb}{}
\lzb{1.}{}{B: z = √(x² + y²) (zweite Option)}{z = √(x² + y²) (zweite Option)}
\end{pfloesungb}
\begin{pfloesungb}{}
\lzb{2.}{}{AB $\approx$ 31,8 m}{AB = √(15² + 28²) = √1009 $\approx$ 31,8 m}
\end{pfloesungb}
\begin{pfloesungb}{}
\lzb{3.}{}{BC $\approx$ 10,9 m}{BC = √(12² $-$ 5²) = √119 $\approx$ 10,9 m}
\end{pfloesungb}
\begin{pfloesungb}{}
\lzb{4.}{}{D: z² = x² + y² (vierte Option)}{z² = x² + y² (vierte Option)}
\end{pfloesungb}
\begin{pfloesungb}{}
\lzb{5.}{}{Radius 2,9 m und Dachhöhe als Katheten, Mantellinie als Hypotenuse: s = √(2,9² + h²)}{Radius 2,9 m und Dachhöhe als Katheten, Mantellinie als Hypotenuse: s = √(2,9² + h²)}
\end{pfloesungb}
\begin{pfloesungb}{}
\lzb{6.}{}{AD $\approx$ 677 m}{AD = √(860² $-$ 530²) = √458700 $\approx$ 677 m}
\end{pfloesungb}
\end{document}